%% ============================================================
%%  Objetivos
%%  Duenio inicial: claude-1
%%  Reglas: ver ../../CLAUDE.md
%%  No editar si TASKS.md marca este archivo IN_PROGRESS por otro agente.
%% ============================================================

\section{Objetivos}

\subsection{Objetivo general}

Desarrollar un framework metodológico agnóstico al dominio que transforme la representación latente aprendida por modelos profundos en una hipótesis relacional compacta, interpretable y evaluable, validada mediante una cadena de evidencia multicriterio que integre estabilidad estadística, fidelidad predictiva y contraste contra paradigmas no probabilísticos.

\subsection{Objetivos específicos}

\begin{enumerate}
\item Formalizar matemáticamente las transformaciones $\Phi$, $\Psi$ y $F$ que componen el pipeline del framework.
\item Evaluar la estabilidad de la representación relacional ante variaciones estocásticas del entrenamiento, distinguiendo el núcleo relacional persistente de la periferia dependiente de la inicialización.
\item Construir una cadena de evidencia acumulativa de cinco eslabones --estabilidad continua, estabilidad estructural, consenso estadístico frente a un modelo nulo explícito, fidelidad predictiva por ablación dirigida y especificidad/robustez--, con umbrales de aprobación e indicadores de fallo declarados antes de la ejecución.
\item Triangular la evidencia mediante cuatro paradigmas ortogonales al enfoque frecuentista --intervencional, geométrico-topológico, de teoría de la información y de baselines deterministas--, de modo que la conclusión sobre una arista no dependa de una única familia de supuestos.
\item Formular una taxonomía diagnóstica que permita distinguir operacionalmente una relación estructural genuina de un atajo predictivo, de un artefacto de tipo sumidero de atención y de ruido estocástico.
\item Validar la hipótesis reconstruida mediante perturbación dirigida sobre el modelo, arista por arista.
\item Demostrar la generalidad del framework mediante su aplicación a un segundo dominio distinto al caso de estudio principal y a un conjunto sintético con estructura conocida.
\end{enumerate}
